\documentclass{article}
\usepackage{amsmath,amsthm,amssymb}
\usepackage {mathtext}
\usepackage{amsfonts} 
\usepackage[T1,T2A]{fontenc}
\usepackage[utf8]{inputenc}
\usepackage [english, russian] {babel}
\usepackage[babel=true]{microtype}
\usepackage{caption}
\usepackage{subcaption}
\captionsetup[table]{position=t,justification=raggedleft,slc=off}
\setlength{\abovecaptionskip}{0pt}
\usepackage {geometry}
\usepackage{graphicx}
\usepackage{multirow}
\geometry{margin=2cm}
\title{Отчёт о лабораторной работе №2.1.4 \\ Определение теплоёмкости твёрдых тел}
\author{Павел Рудачихин \\ группа Б04-507}
\date{\today}

\begin{document}
	\maketitle
	
	\section*{Аннотация}
	
	\hspace{4mm}\textbf{Цель работы:} 1) измерение кривых нагревания $T_{\text{heat}}(t)$ и охлаждения $T_{\text{cool}}(t)$ пустого калориметра и системы «калориметр + твёрдое тело»; 2) определение коэффициента теплоотдачи стенок калориметра; 3) определение теплоёмкости пустого калориметра и удельной теплоёмкости твёрдого тела.
	
	\textbf{В работе используются:} калориметр с нагревателем и термометром сопротивления; универсальный вольтметр В7-78/3 в режиме омметра, измеритель температуры — термопара К-типа совместно с универсальным вольтметром В7-78/2, источник питания GPS-72303, универсальные вольтметры В7-78/3 (в режиме амперметра) и KEITHLEY (в режиме вольтметра) для измерения мощности нагревателя.
	
	\textbf{Методы:} 1) интегральный метод обработки кривых нагревания и охлаждения; 2) дифференциальный метод определения теплоёмкости по производной в «удобной точке»; 3) метод наименьших квадратов для обработки экспериментальных данных.
	
	\section*{Теоретические сведения}

	
	\hspace{4mm} В данной работе измерение теплоёмкости твёрдых тел производится по стандартной схеме. Исследуемое тело помещается в калориметр с нагревателем мощностью $P$. Пусть $\Delta Q$ – количество тепла, подведённое к системе «тело + калориметр» за малое время $\Delta t$, а $\Delta T$ – изменение её температуры. Тогда согласно определению теплоёмкость системы будет равна:
	
	\begin{equation}
		C = \frac{\Delta Q}{\Delta T} \tag{1}
	\end{equation}
	
	В реальных условиях $\Delta Q \neq P\Delta t$, так как часть энергии, выделенной нагревателем, уходит из калориметра из-за теплопроводности его стенок. В результате количество тепла $\Delta Q = C\Delta T$, подведённое к системе, будет меньше $P\Delta t$ на величину тепловых потерь:
	
	\begin{equation}
		C\Delta T = P\Delta t - \lambda(T - T_k)\Delta t \tag{2}
	\end{equation}
	
	где $\lambda$ – коэффициент теплоотдачи стенок калориметра, $T$ – температура тела и калориметра, $T_k$ – температура окружающего калориметр воздуха (комнатная).
	
	\subsection*{Интегральный метод}
	
	При постоянной комнатной температуре ($T_k = \text{const}$) процессы нагревания и охлаждения описываются дифференциальными уравнениями:
	
	\begin{equation}
		C dT_{\text{heat}} = P dt - \lambda [T_{\text{heat}}(t) - T_k] dt \tag{3}
	\end{equation}
	
	\begin{equation}
		C dT_{\text{cool}} = -\lambda [T_{\text{cool}}(t) - T_k] dt \tag{4}
	\end{equation}
	
	Интегрирование уравнения (4) даёт зависимость температуры охлаждения от времени:
	
	\begin{equation}
		T_{\text{cool}}(t) = (T_0 - T_k) e^{-\frac{\lambda}{C}t} + T_k \tag{5}
	\end{equation}
	
	где $T_0$ — температура в момент выключения нагревателя.
	
	Уравнение (5) легко спрямляется в координатах $(\ln(T_{\text{cool}}-T_k), t)$. Тангенс угла наклона этой прямой позволяет определить отношение $\frac{\lambda}{C}$.
	
	Интегрирование уравнения (3) при начальном условии $T_{\text{heat}}(0) = T_k$ даёт:
	
	\begin{equation}
		T_{\text{heat}}(t) = \frac{P}{\lambda}(1 - e^{-\frac{\lambda}{C}t}) + T_k \tag{6}
	\end{equation}
	
	По найденному из кривой охлаждения отношению $\frac{\lambda}{C}$ можно определить $\lambda$, а зная $\lambda$ и $\frac{\lambda}{C}$ — искомую теплоёмкость $C$.
	
	\subsection*{Дифференциальный метод}
	
	При существенных колебаниях комнатной температуры следует использовать дифференциальные методы. К «удобным» точкам относится точка на кривой нагревания, при которой температура калориметра совпадает с комнатной. Дифференцируя уравнение (3) при $T_{\text{heat}}(t) = T_k(t)$, получим:
	
	\begin{equation}
		C = \frac{P}{(dT_{\text{heat}}/dt)_{T=T_k}} \tag{7}
	\end{equation}
	
	Другими «удобными» точками являются точки при одной и той же температуре $T$ на кривых нагревания и охлаждения. В этом случае, решая систему уравнений (3) и (4), получим:
	
	\begin{equation}
		\lambda = \frac{P}{(T - T_{k2})(1 - \frac{A}{B}) + T_{k2} - T_{k1}} \tag{8}
	\end{equation}
	
	\begin{equation}
		C = \frac{P}{A - B + A \frac{T_{k1} - T_{k2}}{T - T_{k1}}} \tag{9}
	\end{equation}
	
	где $A = (dT/dt)_{\text{heat}}$, $B = (dT/dt)_{\text{cool}}$, $T_{k1}$ и $T_{k2}$ – комнатная температура в моменты времени $t_1$ и $t_2$, когда $T_{\text{heat}}(t_1) = T_{\text{cool}}(t_2) = T$.
	
	При равенстве комнатных температур ($T_{k1} = T_{k2} = T_k$) формулы упрощаются:
	
	\begin{equation}
		\lambda = \frac{P}{(T - T_k)(1 - \frac{A}{B})} \tag{10}
	\end{equation}
	
	\begin{equation}
		C = \frac{P}{A - B} \tag{11}
	\end{equation}
	
	\subsection*{Пересчёт сопротивления в температуру}
	
	Температура измеряется термометром сопротивления, сопротивление которого изменяется по закону:
	
	\begin{equation}
		R_T = R_{273}[1 + \alpha(T - 273)] \tag{12}
	\end{equation}
	
	где $\alpha = 4,28 \cdot 10^{-3}$ град$^{-1}$ — температурный коэффициент сопротивления меди.
	
	Выразим $R_{273}$ через измеренное значение $R_k$ — сопротивление термометра при комнатной температуре:
	
	\begin{equation}
		R_{273} = \frac{R_k}{1 + \alpha(T_k - 273)} \tag{13}
	\end{equation}
	
	Подставляя (13) в (12), найдём:
	
	\begin{equation}
		T(R_T) = 273 + \frac{R_T}{\alpha R_k}[1 + \alpha(T_k - 273)] - \frac{1}{\alpha} \tag{14}
	\end{equation}
	
	Для используемой установки формулы пересчёта имеют вид:
	\begin{itemize}
		\item Для установки (1): $T(R_T) = 14.583955001619313455 \cdot R_T + 39.35514018691588785$
		\item Для установки (2): $T(R_T) = 14.377980252039598845 \cdot R_T + 39.35514018691588785$
	\end{itemize}
	
	\section*{Экспериментальная установка}
	
	Установка состоит из калориметра с пенопластовой изоляцией, помещённого в ящик из многослойной клееной фанеры. Внутренние стенки калориметра выполнены из материала с высокой теплопроводностью. Надёжность теплового контакта между телом и стенками обеспечивается их формой: они имеют вид усечённых конусов и плотно прилегают друг к другу.
	
	В стенку калориметра вмонтированы спираль нагревателя (СН) и спираль термометра сопротивления.
	
	\subsection*{Используемое оборудование}
	
	\begin{itemize}
		\item Калориметр с нагревателем и термометром сопротивления
		\item Универсальный вольтметр В7-78/3 (режим омметра) — измерение сопротивления термометра
		\item Универсальный вольтметр В7-78/2 (режим термопары К-типа) — измерение комнатной температуры
		\item Универсальный вольтметр KEITHLEY (режим вольтметра) — измерение напряжения нагревателя
		\item Универсальный вольтметр В7-78/3 (режим амперметра) — измерение тока нагревателя
		\item Источник питания GPS-72303
	\end{itemize}
	
	\subsection*{Параметры эксперимента}
	
	\begin{itemize}
		\item Масса железного образца: $m_{\text{Fe}} = 814.8 \pm 0.5$ г
		\item Масса алюминиевого образца: $m_{\text{Al}} = 294.1 \pm 0.5$ г
		\item Показание вольтметра (напряжение): $U = 27.3408$ В
		\item Показание амперметра (ток): $I = 0.229247$ А
	\end{itemize}
	
	Мощность нагревателя рассчитывается по формуле $P = UI$:
	
	\begin{equation}
		P = 27.3408 \cdot 0.229247 = 6.268 \text{ Вт}
	\end{equation}


\section*{Заключение}

\hspace{4mm} Значения коэффициентов Джоуля-Томсона согласуются со справочными данными. Однако значения констант уравнения Ван-дер-Ваальса завышены примерно в 4 раза. Вместе с тем температура инверсии совпадает со справочной в пределах погрешности, так как зависит от отношения констант, которое соответствует табличному. 

Причиной завышения констант $a$ и $b$ могли быть:

\begin{itemize}
	\item крайне малое количество экспериментальных данных, ибо по трём значениям невозможно установить зависимость с необходимой точностью;
	\item близость к критической температуре, что влияет на свойства газа;
	\item грубые приближения, например, разрежённость газа или теплоизоляция трубки.
\end{itemize}

Так или иначе, цели работы достигнуты: определено изменение температуры углекислого газа при протекании через малопроницаемую перегородку при разных начальных значениях давления и температуры; 2)по результатам опытов вычислены коэффициенты Джоуля-Томсона, константы a и b в модели Ван-дер-Ваальса.

\end {document}